\section{Introduction}

Household inflation expectations matter for policy transmission because communication, wage bargaining, and precautionary spending all react to what households believe about near-term prices. The Chinese case is empirically demanding: uncertainty shocks, geopolitical episodes, and policy communication shifts often arrive together, while the official CPI process includes sizable food-driven volatility. A design that confounds these channels can still fit the data and yet fail as policy guidance, so I center the analysis on testable prediction chains rather than narrative correlations that would collapse under referee scrutiny.

Existing evidence for China documents expectation errors but does not cleanly separate measurement from identification, and most dynamic exercises rely on models where the shock definition and the result are mechanically linked. That leaves a gap exactly where policy relevance is largest, because a policymaker needs to know whether revisions predict subsequent errors after accounting for uncertainty states and whether this relation survives falsification tests instead of disappearing when controls change. The missing piece is a design where salience variation has an external source and where competing mechanisms can be rejected in the same empirical frame.

I address that gap by combining three layers. I quantify survey-based expectations with a Carlson--Parkin mapping that keeps uncertain responses as information through bound analysis, then estimate a salience-IV design in which media congestion instruments inflation salience and traces the chain from salience to revisions to forecast errors. I map dynamic responses with local projections, and I keep Bayesian modules visible as diagnostics rather than as the causal spine: a state-space model tracks the time variation of the diagnostic coefficient with full posterior workflow checks, while an appendix BVAR supplies structural consistency checks anchored to a salience proxy instead of circular sign logic.

The estimates show a negative revision--error slope in baseline and IV specifications, with magnitude, standard errors, and sample sizes reported in Tables~\ref{tab:identification_main} and \ref{tab:ols_baseline}. Placebo outcomes, permutation instruments, and anticipation tests remain weak relative to the main effect in Table~\ref{tab:placebo_tests}, and local projections in Table~\ref{tab:lp_dynamic} show fast error correction with a smaller inflation response over the same horizons. A policy rule that adjusts raw survey expectations by the estimated diagnostic component lowers forecast loss in backtests reported in Table~\ref{tab:policy_backtest}, which moves the paper from behavioral interpretation to operational policy mapping where each forecast gain links to estimated coefficients and uncertainty states.
